%% kapa: exam=1 round=34 year=2023 period=1 form=A subject=civil_law subject_ko=민법 seq=40 number=40
%% source: Q-Net 감정평가사 기출문제 (cst00309) artlSeq=5211789 file=raw/1008/5211789/2023년도 제34회 감정평가사 1차시험 1교시.pdf
%% conversion: kordoc 4.24.51 -> tools/gen1.py (auto); review=proofread-sonnet; flags=none
\begin{question}{40}
법정지상권의 성립에 관한 설명으로 옳지 않은 것은? (다툼이 있으면 판례에 따름)
\begin{choices}
  \item 토지에 저당권이 설정된 후에 저당권자의 동의를 얻어 건물이 신축된 경우라도 법정지상권은 성립한다.
  \item 토지의 정착물로 볼 수 없는 가설 건축물의 소유를 위한 법정지상권은 성립하지 않는다.
  \item 무허가건물이나 미등기건물을 위해서도 관습법상의 법정지상권이 인정될 수 있다.
  \item 토지공유자 중 1인이 다른 공유자의 동의를 얻어 그 지상에 건물을 소유하면서 자신의 토지지분에 저당권을 설정한 후 그 실행경매로 인하여 그 공유지분권자와 건물소유자가 달라진 경우에는 법정지상권이 성립하지 않는다.
  \item 동일인 소유의 토지와 건물 중 건물에 전세권이 설정된 후 토지소유자가 바뀐 경우, 건물소유자가 그 토지에 대하여 지상권을 취득한 것으로 본다.
\end{choices}
\end{question}
\answer{1}
